\textbf{Budget and tuning.} The registered 400-step budget was too short for the GRU; the verdicts changed at 1000 and 2000 steps and may change again beyond 2000. The budget sweep is post hoc; it covers all horizons only for regime and etth1 and $H=96$ for the other four tasks, and the noise sweep and design ablation exist at 400 steps only. Learning rates and model sizes were fixed in advance and never tuned, so a tuned model of either family could change the ranking. \textbf{Seeds.} Three seeds per cell make the win rule a coarse screen: sign-consistent effects below 5\% (PatchTST on regime at 400 steps) are reported as sub-threshold, and there are no formal significance tests. \textbf{Synthetic design.} All synthetic conclusions depend on our own generators (sinusoids, a three-state regime process that rescales periods and amplitudes, mild trends, AR(1) noise). Other nonlinearities (state-dependent dynamics, cross-channel interactions, heavy tails) are absent, and channels and models are independent per channel. \textbf{Reimplemented baselines.} The patch transformer and GRU are small reimplementations, not the published configurations; a result here says nothing about the original PatchTST. \textbf{One real dataset.} ETTh1 is a single series; seeds vary only initialisation and batch order. \textbf{Confounds.} Instance normalisation is used for all learned models and changes MSE on regime and trend by amounts comparable to the model gaps (Table~\ref{tab:abl}), so the ranking could depend on it. \textbf{Status of the sweeps.} The dwell and noise sweeps were planned in the brief as the tests of H2 and H3, but the win checks on individual sweep cells are not registered hypotheses. A full study would add seeds, tuned baselines, budgets to convergence in every cell, more real datasets and crossed sweeps.
